\chapter{Thermal ensembles and response functions}
\label{narr:thermal}

At finite temperature the choice of ensemble is as important as the choice of
Hamiltonian. A fixed flux-sector matter calculation, a sum over physical flux
sectors, and an unprojected four-flavor Gaussian ensemble count different states.
A two-peak curve alone cannot identify which of these has been computed.

% EDITORIAL QUESTION: What thermodynamic identities hold for a fixed physical Hamiltonian?
\section{Finite physical Gibbs states}
For a fixed $H$ on $\mathcal H_s$, $T>0$, and $\beta=1/T$, define
\begin{equation}
 Z=\Tr_{\mathcal H_s}e^{-\beta H},\quad
 \rho_\beta=Z^{-1}e^{-\beta H},\quad
 F=-T\ln Z,\quad E=\Tr(\rho_\beta H).
\end{equation}
The finite Gibbs identities are
\begin{equation}
 S=-\Tr\rho_\beta\ln\rho_\beta=\beta E+\ln Z,
 \qquad C_V=\frac{\dd E}{\dd T}=\beta^2(\langle H^2\rangle-E^2)\geq0.
 \label{eq:narr:gibbs}
\end{equation}
At infinite temperature $S/N\to\ln2$ for $N$ spins. The maximally mixed
four-Majorana auxiliary space instead has $S/N=\ln4$ before projection.
An entropy plateau therefore needs a declared state count and normalization,
not only a comparison to a familiar fraction of $\ln2$.

For a \emph{fixed} quadratic auxiliary matrix $A$, the independent-mode thermal
state of Eq.~\eqref{eq:narr:one_mode} has
\begin{equation}
 \Omega_T=\ii\tanh(\ii\beta A/2).
 \label{eq:narr:thermal_covariance}
\end{equation}
At a zero eigenvalue the finite-temperature occupation is one half, not a
specified pure-state choice. In a self-consistent approximation $A=A(T)$, the
energy functional also contains interaction/double-counting terms. Its total
$T$ derivative cannot be replaced by the variance formula for a fixed quadratic
$A$. A thermodynamically consistent variational free energy and its stationary
derivatives must be derived for that approximation.

% EDITORIAL QUESTION: What can a binary plaquette toy establish about a flux crossover?
\section{A diagnostic toy and the missing ensemble}
An isolated binary variable with energies $0$ for $W=+1$ and $\Delta_v>0$ for
$W=-1$ has partition function $1+e^{-\beta\Delta_v}$ and
\begin{equation}
 \langle W\rangle=\tanh\frac{\Delta_v}{2T},\qquad
 \rho_-=(1+e^{\beta\Delta_v})^{-1}.
 \label{eq:narr:binary_flux}
\end{equation}
This smooth crossover is an exact result for that toy. With $\Delta_v=0.137$
it gives approximately $0.9979$ at $T=0.02$ and $0.4274$ at $T=0.15$.
It does not give a 90\% loss across that window. On a physical torus the flux
constraints, matter spectrum, and parity restrictions couple the weights;
Eq.~\eqref{eq:narr:binary_flux} cannot be adopted as their derivation.

For any finite Hilbert space and real $T>0$, $Z$ is a positive finite sum of
exponentials, so $F(T)$ is analytic. A true first-order nonanalyticity concerns
a controlled thermodynamic limit. A finite-system model comparison can assess
which finite curves describe the data, but a large BIC difference cannot supply
that limit or repair a wrong ensemble. This is a prerequisite for reopening the
proposed thermal Wilson-loop task.

% EDITORIAL QUESTION: Which response normalization and selection rule define a probe threshold?
\section{Lehmann weights, detailed balance, and observables}
For a Hermitian physical probe $A$, use the Fourier convention
$S_{AA}(\omega)=\int_{\RR}\dd t\,e^{\ii\omega t}
\Tr(\rho_\beta A(t)A)$, with $A(t)=e^{\ii Ht}Ae^{-\ii Ht}$. Then
\begin{equation}
 S_{AA}(\omega)=\frac{2\pi}{Z}\sum_{m,n}e^{-\beta E_m}|A_{mn}|^2
  \delta\bigl(\omega-(E_n-E_m)\bigr).
 \label{eq:narr:lehmann}
\end{equation}
Completeness and exchanging $m,n$ give
\begin{equation}
 \int\frac{\dd\omega}{2\pi}S_{AA}(\omega)=\langle A^2\rangle,
 \qquad S_{AA}(-\omega)=e^{-\beta\omega}S_{AA}(\omega).
 \label{eq:narr:kms}
\end{equation}
For a Pauli spin component the full integrated weight is one; for $S^\alpha$
it is one quarter. A broadened, truncated frequency window need not integrate
to that value unless its omitted tails are accounted for. Conversely, imposing
the sum rule by normalization does not verify the positions or physical origin
of spectral peaks.

A local spectral peak becomes an STS prediction only after specifying the
probe coupling, tunneling kernel, and experimental observable. Effective
three-spin-coupling spectra and finite-field spin spectra also need their
Hamiltonian bridge. The historical dynamics appendix preserves exploratory
calculations and reproduction language with explicit review warnings. This
book does not present them as a completed reproduction of arXiv:2209.12930.

% EDITORIAL QUESTION: How will this narrative grow beyond the current model scope?
\section{Next questions without inherited certification}
The foundation now supplies tests that a new result must satisfy: the same
physical state space, an explicit convention map, the correct variational or
thermal functional, and a matched observable. Future chapters can extend the
bulk and defect physics as those results are reviewed. New model families can
become separate parts sharing the finite-state foundations. Chapters are
therefore destinations for an evolving argument, while claim identifiers and
source records remain stable across reordering.

The immediate unresolved work includes physical-sector resolvent thresholds,
the sign convention in the legacy Zeeman generator, the exact versus mean-field
prefactor comparison, and the braiding transport test. None is settled by
assembling a book. Their current status and evidence remain in the wave audit
and acceptance register; the appendices retain the route needed to investigate
them without discarding earlier findings.
